\scanpage{38}
% archive foliation: 37

Je me reporte au diagramme (23)
\[
\diagram{\begin{tikzcd}[column sep=large]
&&1\arrow[d]&&\\
&&\mathbb{F}_2^{\underline\omega}\arrow[dl]\arrow[dr,equal]&&1\arrow[dl]\\
&GL_\Lambda(M)\arrow[dl]\arrow[dr]&&\mathbb{F}_2^{\underline\omega}\simeq V(S_Q)\arrow[dl]&\\
\unsure{\mathcal{G}_I|\mathfrak S_I}\arrow[dr,"\inserted{\mathrm{sg}}"']&
\text{cart}&\mathfrak D_Q\arrow[dl,"\chi_C"]&\\
&\{\pm1\}&&
\end{tikzcd}}
\]
qui exprime $GL_\Lambda(M)$ comme produit fibré, la donnée de $M$ équivalant à celle du triple $(I,Q,k)$, où $I$ est un ens. à quatre éléments, $Q$ un carré \inserted{combinatoire}, et $k$ un isom
\[
k:\underset{\text{codiagonales}}{\underbrace{C_Q}}
\ \simeq\ \underline\omega_I.
\]

En termes de $Q$, on a
\[
C_\Delta\simeq F(\Delta)/M_0\simeq
P_{M,\mathbb{F}_2^{\underline\omega}}/\mathbb{F}_2
=S_Q/\mathbb{F}_2\simeq D_Q
\]
ens des diagonales de $Q$ — donc la donnée de l'octaèdre $\Delta$ (au lieu de $M\longleftrightarrow(I,Q,k)$) équivaut à celle de $I$, et des deux ens. à deux éléments $C$ $(\simeq\underline\omega_I)$ et $D$ correspondant aux diagonales et codiagonales du carré $Q$, \deleted{\unsure{dont}} la donnée équivaut à celle d'un torseur sous
$\deleted{\unsure{\mathfrak D_Q}}\mathfrak D_4/Z(\mathfrak D_4)\simeq\deleted{\unsure{\mathbb{F}_2}}\mathbb{F}_2\times\mathbb{F}_2$,
qu'il faut donc relever en un torseur sous $\mathfrak D_4$. Ou encore, décrire la donnée d'un carré combinatoire, dont on a déjà les deux ensembles \inserted{$C$ et $D$} (des \uline{codiagonales} et \uline{diagonales}\ldots\quad$\big]\big]$
